\subsection{Whole Number Exponents} 
Suppose we have a power function of the form
\[
f(x)=a(x-h)^n-k
\]
where $h,k$ are \makeblank[1in] numbers and $n$ is a \makeblank[1in] number. (Recall the \makeblank[1in] are the numbers \makeblank. We ignore $n=\makeblank[.5in]$, since then the function is \makeblank[1in].)
\subsubsection*{Domain}
If we were to multiply out $(x-h)^k$, it is a \makeblank[1.5in], so its domain is \makeblank.
\subsubsection*{Range}
\begin{itemize}
\item If $n$ is \makeblank[1in], $(x-h)^n$ can be \makeblank[1.5in], so the range of $f$ is \makeblank.
\item If $n$ is \makeblank[1in], $(x-h)^n$ must be \makeblank[1.5in], so the range of $f$ depends on $a$ and $k$.
\begin{itemize}
\item If $a$ is \makeblank[1in], the range of $f$ is \makeblank.
\item If $a$ is \makeblank[1in], the range of $f$ is \makeblank.
\end{itemize}
\end{itemize}
\subsubsection*{Inverse}
\begin{itemize}
\item if $n$ is \makeblank[1in], $x^n$ has an inverse, \makeblank.
\item if $n$ is \makeblank[1in], $x^n$ is not \makeblank, but we can fake partial inverses:
\begin{itemize}
\item If we restrict the domain to \makeblank[1in], the inverse is \makeblank[1in]
\item If we restrict the domain to \makeblank[1in], the inverse is \makeblank[1in]
\end{itemize}
\end{itemize}